\subsection{The quotient question in its written generality}

\paragraph{Theorem E4.Q (refutation of 57:309(b)).}
The manuscript asks for ``existence du pr\'esch\'ema quotient''
$\underline{\mathrm{Pic}}_{X/S}/\underline{\mathrm{Pic}}^\tau_{X/S}$.
In this generality the answer is negative, even when the tau-part is flat.
Let
\[
 S=\operatorname{Spec}\mathbf R[[t]],\qquad
 X=V(x^2+y^2-tz^2)\subset\mathbf P^2_S .                  \tag{E4.Q.1}
\]
This is Mumford's example as proved in
\cite[Example~4.14, pp.~31--33]{Kleiman2005}. Its Picard fppf sheaf
$P$ is an algebraic space, not a scheme. After the finite \'etale base change
$S'=\operatorname{Spec}\mathbf C[[t]]$, its component of total degree $n$
is obtained by gluing copies $S'_{a,b}$, $a+b=n$, along their generic point.
Complex conjugation interchanges the origins $(a,b)$ and $(b,a)$.
For $a\ne b$ these origins have no common affine neighborhood: their two
trait sections agree generically, contrary to separatedness of such an
affine neighborhood. If $P$ were a scheme, an affine neighborhood of their
common descended point would pull back to just such an affine neighborhood.

In this example $P^\tau$ is exactly the identity subgroup $S$. Indeed the
geometric generic conic has Picard group $\mathbf Z$, and the geometric
special fibre, two projective lines meeting in one point, has Picard group
$\mathbf Z^2$ by multidegree. Both have trivial neutral and torsion parts.
In the displayed scheme $P_{S'}$, the tau-subfunctor is consequently the
identity copy $S'_{0,0}$; the explicit representation also rules out any
additional infinitesimal structure. This identification descends to $S$.
Therefore, as fppf sheaves of groups,
\[
 P/P^\tau=P/S\simeq P,
\]
which is not a scheme. Here $S$ denotes the trivial group object, not a
nontrivial quotient of the underlying base. Its identity immersion into
$P$ is open and not closed. Thus this example is consistent with the flat-only algebraic-space assertion
in the prompt summary of H6. The stronger scheme conclusion in the fuller
\texttt{HUNCHES.md} formulation additionally assumes closedness, which fails here.

The stronger variant that presupposes $P$ itself schematic, or imposes
closedness of $P^\tau$ while dropping flatness, is a distinct question.
No answer to those extra variants is inferred from (E4.Q.1). In particular,
this example has a geometrically reducible special fibre and gives no
counterexample to the geometrically integral or normal question on 57:315.

\paragraph{Two further consequences of E4.H.}
The character group in (E4.H.1) is killed by two. Its identity section is
not open: in the component $\mu_{2,R}^{\,2}$, the four geometric generic
points specialize to the same special identity. Thus the image of the
whole group under $[2]$ is not open. Question 57:310(4) is therefore false
even for a smooth projective family, extending the simpler cuspidal
field example.

For the margin question about ``Multiplicit\'e de
$\underline{\mathrm{Pic}}^\tau$'' and specialization on 57:310, the same
example makes several precise natural interpretations explicit. The
generic fibre has four reduced points, each of multiplicity one. The
special fibre consists of two copies of $\mu_{2,k}^{\,2}$, each of length
four. Total length changes from four to eight, the number of geometric
components from four to two, and the length of the neutral geometric
component from one to four. Thus none of these numerical multiplicities
is specialization invariant without further hypotheses. The manuscript's
additional words about existence and descent do not specify an object;
we do not replace them by a different theorem.
